\clearpage
\chapter{Conclusiones}

\section{Cumplimiento de objetivos}

El objetivo general consistía en diseñar y evaluar una metodología para generar escenarios sintéticos de retornos y utilizarlos en una optimización de carteras. Este objetivo se ha materializado en un pipeline que comienza con precios ajustados, construye una matriz de 2.766 log-retornos, entrena un VAE, genera 5.000 escenarios, estima parámetros y evalúa pesos sobre 554 observaciones posteriores. La comparación con tres referencias permite situar el resultado frente a reglas simples e históricas.

Los objetivos teóricos se abordaron mediante la revisión de Markowitz, riesgo, hechos estilizados, teoría de la información y modelos generativos. El análisis distingue la función de la divergencia KL en el VAE de la función de máximo Sharpe de la cartera. También diferencia el estado del arte GAN/WGAN de la implementación real, una precisión necesaria para que el alcance sea defendible.

Los objetivos relacionados con datos se cumplieron mediante adquisición reproducible, tratamiento de ausencias, retornos, EDA y partición temporal. La matriz final no contiene nulos y las figuras documentan colas, volatilidad variable y correlaciones. La preparación ajusta el escalador solo con entrenamiento, reduciendo fuga de información.

En modelado, además del VAE principal, se ejecutó una comparación real con AE determinista y VAE de latente reducido bajo tres semillas. Esta ampliación responde a la necesidad de justificar el modelo y muestra que reconstrucción y generación son problemas relacionados, pero no equivalentes. La evaluación financiera completa el ciclo utilizando métricas comunes.

\section{Principales resultados obtenidos}

El VAE principal obtuvo MSE 0,00777 en entrenamiento y 0,01049 en test; el MAE pasó de 0,06022 a 0,06857. La brecha es moderada, pero confirma que la capacidad de generalización no es perfecta. Los resúmenes sintéticos muestran que parte de la volatilidad se suaviza, un aspecto crítico cuando las muestras se usan para estimar riesgo.

En la comparación de 120 épocas, el AE logró test MSE 0,006308 frente a 0,010425 del VAE-8 y 0,010561 del VAE-4. También presentó menores errores de volatilidad y correlación bajo el esquema empírico de muestreo. El VAE-8, en cambio, obtuvo el menor error de medias sintéticas (0,003852) y ofrece un prior explícito. Reducir el latente a cuatro dimensiones no mejoró el conjunto de métricas.

La cartera VAE alcanzó retorno acumulado 242,32~\%, retorno anualizado 75,02~\%, volatilidad 25,08~\%, Sharpe 2,99 y máximo \textit{drawdown} -13,96~\%. En ese mismo test, Markowitz restringido obtuvo 101,38~\%, 17,01~\% de volatilidad y Sharpe 2,20; la equiponderada presentó la menor volatilidad y caída. El VAE ofrece, por tanto, la mayor rentabilidad ajustada por volatilidad de la tabla, pero también el mayor riesgo absoluto.

\section{Valoración final de la metodología}

CRISP-DM resultó útil para conectar decisiones que podrían haberse presentado de manera fragmentada. La comprensión del problema fijó benchmarks y cautelas; la comprensión de datos produjo EDA verificable; la preparación evitó mezclar futuro y pasado; el modelado conservó configuraciones; y la evaluación empleó un test común. La reproducibilidad se apoya en scripts, semillas y artefactos intermedios.

La separación temporal es una fortaleza frente a una evaluación sobre entrenamiento. No obstante, un único corte sigue siendo débil para mercados no estacionarios. La selección del modelo y de restricciones debería realizarse dentro de ventanas anteriores, reservando múltiples periodos para evaluación. La metodología actual debe entenderse como primera prueba controlada.

La comparación multimodelo añade una conclusión metodológica: escoger por reconstrucción habría favorecido al AE, mientras el propósito generativo favorece valorar regularidad y momentos. Ninguna métrica resume por sí sola la idoneidad. En futuros trabajos conviene definir un criterio compuesto antes de observar el resultado financiero.

\section{Aportaciones del trabajo}

La primera aportación es un pipeline integral y auditable que conecta generación y optimización sin ocultar la capa clásica final. La segunda es un conjunto de productos reproducibles: EDA, modelos, escenarios, pesos y métricas. La tercera es una lectura crítica que evita presentar GAN/WGAN como ejecutadas y sitúa KL en su función correcta.

La cuarta aportación es empírica: bajo un mismo test, la fuente sintética modifica sustancialmente la cartera y produce métricas diferentes a la estimación histórica. Este hallazgo demuestra sensibilidad al procedimiento de estimación, no superioridad universal. La quinta es la comparación AE/VAE, que identifica oportunidades concretas para mejorar la preservación de volatilidades y correlaciones.

\section{Conclusión general}

La principal aportación del trabajo no reside en demostrar que un VAE sea universalmente superior a los métodos clásicos, sino en documentar un procedimiento reproducible que permite conectar modelos generativos, generación de escenarios y optimización interpretable de carteras bajo un protocolo evaluable.

Los resultados sugieren que un VAE puede utilizarse de forma técnicamente viable para generar escenarios multivariantes y alimentar una optimización interpretable. En el periodo analizado, la cartera resultante obtiene métricas favorables de retorno y Sharpe. Sin embargo, la magnitud del resultado, la mayor volatilidad y las discrepancias estadísticas de las muestras exigen cautela.

No puede afirmarse que el modelo prediga el mercado, garantice rendimientos o sustituya una validación financiera completa. El valor del TFM reside en demostrar y documentar el encadenamiento experimental, comparar alternativas y hacer visibles sus supuestos. La contribución constituye una base sólida para ampliar ventanas, modelos y restricciones, pero no una recomendación de inversión.
